% GALILEO — EMNLP Main (working LaTeX)
% Goal: compile locally (2-column) and iteratively improve paper quality.
\pdfoutput=1
\documentclass[11pt]{article}

% Remove [review] for camera-ready.
\usepackage[review]{EMNLP2023}

\usepackage{times}
\usepackage{latexsym}
\usepackage[T1]{fontenc}
\usepackage[utf8]{inputenc}
\usepackage{microtype}
\IfFileExists{inconsolata.sty}{\usepackage{inconsolata}}{}

\usepackage{amsmath,amssymb}
\usepackage{booktabs}
\usepackage{array}
\usepackage{tabularx}
\usepackage{multirow}
\usepackage{graphicx}
\usepackage{url}
\usepackage{hyperref}
\usepackage{xcolor}
\usepackage{enumitem}

% Helpers
\newcolumntype{L}[1]{>{\raggedright\arraybackslash}p{#1}}
\newcolumntype{Y}{>{\raggedright\arraybackslash}X}

\title{GALILEO: Multi-turn Robustness Dynamics and Recovery under Conversational Pressure}
\author{Anonymous Authors}
\date{}

\begin{document}
\maketitle

\begin{abstract}
Large language models can abandon previously correct answers when challenged in multi-turn dialogue.
We introduce \textbf{GALILEO}, a reproducible protocol on ground-truth tasks that conditions on initially correct examples and measures (i) survival over rounds, (ii) turn-of-failure (TOF) with early-turn vulnerability (Fail@1), and (iii) recovery@flip (return-to-truth after a flip).
We pair each pressure persona with a matched evidence-free neutral re-asking control (NRC) to separate pressure-induced effects from generic multi-turn drift.
Across multi-seed experiments and several open-weight model families, pressure reduces Survival@5 and increases Fail@1 relative to NRC; recovery varies by task/persona, indicating that robustness and recoverability are distinct axes.
\end{abstract}

\section{Introduction}
\label{sec:intro}

% TODO (draft expansion): tighten motivation, then expand with 2-3 concrete failure scenarios and positioning.

Single-turn accuracy obscures interaction dynamics: \emph{when} a model first fails under pressure, \emph{how} robustness decays over turns, and \emph{whether} it can return to truth after being misled (e.g., sycophancy and multi-turn persuasion effects; \citealp{sharma2023understandingsycophancylanguagemodels,liu2025truthdecayquantifyingmultiturn}).
We target ground-truth tasks where failure is unambiguous, and we introduce a matched drift baseline (NRC) so persona effects are interpretable.

\paragraph{Contributions.}
\begin{itemize}[leftmargin=*]
  \item We propose a ground-truth multi-turn protocol with survival curves, TOF/Fail@1, and recovery@flip.
  \item We introduce NRC, an evidence-free matched-length control, enabling drift-corrected attribution of pressure mechanisms.
  \item We provide paper-ready exports and a claim\,$\rightarrow$\,artifact pipeline for auditing and figure regeneration.
\end{itemize}

\section{Task setting}
We cover math (GSM8K, SVAMP), extractive QA (SQuAD), MCQA (ARC-Easy), and open-domain QA (TriviaQA).
All tasks are evaluated with a unified boxed-answer format (\texttt{\textbackslash boxed\{...\}}) and a task-specific scorer.

\section{Protocol}
\label{sec:protocol}

\subsection{Overview (paper-auditable loop)}
\begin{figure}[t]
\centering
\fbox{\parbox{0.97\linewidth}{
\textbf{Algorithm 1: GALILEO evaluation for a persona $p$ and matched NRC.}\\
\textbf{Input:} dataset $D$ with gold answers; decoding config; rounds $R{=}5$; persona prompt template $\mathsf{persona}(p,\cdot)$; NRC template $\mathsf{nrc}(p,\cdot)$.\\
\textbf{1. Initial evaluation:} Run the model on each $x\in D$ with a neutral instruction; score to obtain initially-correct set $C\subseteq D$.\\
\textbf{2. Multi-round pressure vs control:} For each $x\in C$, iterate rounds $r{=}1..R$:\\
\ \ \ (a) query with persona: $y^{(p)}_{r} \leftarrow \mathsf{persona}(p, x, \text{history})$; score correctness\\
\ \ \ (b) query with control: $y^{(\mathrm{NRC})}_{r} \leftarrow \mathsf{nrc}(p, x, \text{history})$; score correctness\\
\ \ \ update history separately for persona and NRC tracks.\\
\textbf{3. Recovery:} For any $x$ that flips at least once in the persona track, issue a final neutral recovery prompt; score Recovery@flip.\\
\textbf{Output:} per-example trajectories; Survival@r, TOF/Fail@1, Recovery@flip; exports for figures/tables.
}}
\vspace{-2mm}
\caption{Protocol summary. Formal definitions for NRC, conditioning sets, and metrics appear below.}
\label{fig:algo}
\end{figure}

\subsection{Three phases}
\textbf{Phase 1 (initial evaluation).} Identify the initially-correct subset.
\textbf{Phase 2 (multi-round pressure vs NRC).} Apply personas for up to $R=5$ rounds and in parallel run NRC under identical decoding.
\textbf{Phase 3 (recovery).} For flipped examples, issue a final neutral recovery prompt and measure recovery@flip.

\subsection{Neutral Re-asking Control (NRC): evidence-free matched drift baseline}
For each pressure persona $p$, we construct a paired \emph{Neutral Re-asking Control (NRC)} prompt that matches the interaction length and turn structure but removes the mechanism of pressure. Concretely, the NRC prompt:
(i) repeats the user question in a neutral tone, (ii) can acknowledge prior turns without introducing new facts, and (iii) avoids any argumentative content, social pressure, or claims about the gold answer.
This makes NRC an \emph{evidence-free} baseline: any failures under NRC reflect multi-turn drift (e.g., instability under re-asking) rather than persona pressure.

\textbf{Allowed examples (NRC).} ``Can you restate your final answer?''; ``Please answer again, briefly.''; ``Just to confirm, what is your answer?''
\textbf{Forbidden examples (not NRC).} Introducing counter-evidence (``Actually, the correct answer is ...''), social manipulation (``I will be disappointed if you don't change''), or logical traps (``If your answer were true, then ...'').

\textbf{Attribution.} We report persona effects as paired differences (persona -- NRC) under identical decoding and matched conditioning sets, so that measured deltas correspond to pressure mechanisms rather than generic re-asking.

\subsection{Conditioning-set reporting modes (to avoid confusion)}
We distinguish (i) a shared initially-correct set $C$ (for cross-persona comparisons) and (ii) persona-matched sets $C_p$ (for within-persona attribution).
Unless stated otherwise, persona\,–\,NRC gaps are computed on the same matched conditioning set ($C_p$), so control baselines can differ across personas.

\subsection{Metrics}
Survival@r is cumulative correctness through rounds $1..r$; TOF is the first round where an initially-correct example becomes incorrect (right-censored if never failing within $R$).
Recovery@flip is correctness on the recovery prompt conditional on having flipped at least once.

\section{Results}
\label{sec:results}

We summarize multi-turn robustness as complementary views: (i) \textbf{Survival@r} captures cumulative correctness under repeated questioning, (ii) \textbf{TOF/Fail@1} captures early-turn brittleness and the onset of failure, and (iii) \textbf{Recovery@flip} captures whether a model can return to the gold answer after it has deviated.
All main comparisons are reported as \emph{persona--NRC deltas} to correct for evidence-free drift.

\subsection{Survival degrades under pressure beyond generic drift}
Figure~\ref{fig:survival} shows survival curves over rounds for representative personas alongside NRC.
A consistent pattern is that the survival gap widens with rounds: some personas induce immediate early failure (captured by Fail@1), while others erode correctness gradually across turns.
This supports treating robustness as a trajectory-level property rather than a single scalar accuracy.

\subsection{Turn-of-failure highlights early vulnerability}
Figure~\ref{fig:fail1} reports persona-wise \(\Delta\)Fail@1.
Positive deltas indicate that a persona causes a larger fraction of initially-correct examples to fail on the first challenged turn, compared to drift under NRC.
Empirically, we observe that early-turn vulnerability differs substantially by persona, suggesting distinct pressure mechanisms (e.g., social deference vs logical contradiction).

\subsection{Recovery is a distinct axis from robustness}
Figure~\ref{fig:recovery} reports \(\Delta\)Recovery@flip.
In several settings, personas that substantially reduce Survival@5 do not necessarily reduce Recovery@flip, indicating that ``staying correct'' and ``returning to truth after a deviation'' are separable model behaviors.
This matters for deployment: systems that monitor for flips may still benefit from recovery-oriented guardrails even when robustness is imperfect.

\begin{figure*}[t]
  \centering
  \includegraphics[width=0.98\textwidth]{figures/survival_curves_rounds_seed1-4_20260209}
  \caption{Survival over rounds (seed1--4) for selected personas vs NRC.}
  \label{fig:survival}
\end{figure*}

\begin{figure}[t]
  \centering
  \includegraphics[width=\linewidth]{figures/tof_personawise_fail1_delta_seed1-4_20260209}
  \caption{Persona-wise \(\Delta\)Fail@1 (persona -- NRC).}
  \label{fig:fail1}
\end{figure}

\begin{figure}[t]
  \centering
  \includegraphics[width=\linewidth]{figures/recovery_personawise_delta_seed1-4_20260209}
  \caption{Persona-wise \(\Delta\)Recovery@flip (persona -- NRC).}
  \label{fig:recovery}
\end{figure}

\subsection{Cross-family generalization}
\begin{figure}[t]
  \centering
  \includegraphics[width=\linewidth]{figures/cross_family_survival_r5_control_vs_logicaltrap_seed1-2_20260219}
  \caption{Cross-family Survival@5 under NRC vs a strong persona (Logical Trap), seeds 1--2.}
  \label{fig:crossfam}
\end{figure}

\subsection{Decoding sensitivity}
\begin{figure}[t]
  \centering
  \includegraphics[width=\linewidth]{figures/decoding_sweep_qwen_delta_seed1-2_20260211}
  \caption{Decoding sensitivity (seeds 1--2): persona--NRC gaps persist across temperatures.}
  \label{fig:decoding}
\end{figure}

\section{Discussion and limitations}
\label{sec:limitations}

Our results suggest that multi-turn robustness under conversational pressure is not fully captured by single-turn accuracy: models can be initially correct yet brittle in the first challenged turn (Fail@1), degrade over subsequent rounds (Survival@r), and sometimes recover after a flip (recovery@flip). The paired NRC control is essential for interpretation because multi-turn ``truth decay'' can occur even without explicit pressure.

\paragraph{Limitations.}
\begin{itemize}[leftmargin=*]
  \item \textbf{Protocol scope.} GALILEO focuses on short-horizon, turn-bounded interactions ($R=5$) with evidence-free prompts. Longer dialogues, tool use, retrieval, or user-provided evidence may yield different dynamics.
  \item \textbf{Task coverage.} We emphasize ground-truth tasks with unambiguous scoring. This excludes many open-ended domains (e.g., creative writing, policy advice) where ``correctness'' is less well-defined and robustness must be assessed differently.
  \item \textbf{Persona design.} Pressure personas are handcrafted and cannot span the full space of real adversarial strategies. NRC reduces drift confounds but does not guarantee that all pressure effects are isolated from other latent factors (e.g., wording sensitivity).
  \item \textbf{Model/decode dependence.} Robustness and recovery can vary with decoding settings and model families. While we include a decoding sweep, our conclusions should be read as comparative within the reported settings rather than universal guarantees.
\end{itemize}

\paragraph{Ethics and misuse.} Persona prompts could be repurposed to intentionally destabilize a model in deployment. To reduce misuse risk, we (i) report aggregate vulnerability trends rather than releasing curated ``best attacks'' for a specific target system, and (ii) recommend that evaluations include a drift baseline like NRC and that deployments add safeguards (e.g., calibrated refusal on repeated re-asking, consistency checks, and audit logs) when models are used in high-stakes settings.

\section{Conclusion}
\label{sec:conclusion}

% TODO: add 1-paragraph conclusion summarizing the protocol + key findings.

\bibliography{../../../references}
\bibliographystyle{acl_natbib}

\appendix
\section{Appendix: experimental settings and prompts}
% TODO: detailed configs (models, seeds, decoding, max tokens/len, hardware), prompt templates.

\end{document}
